\begin{table}[p]\centering
\def\sym#1{\ifmmode^{#1}\else\(^{#1}\)\fi}
\begingroup
\renewcommand{\arraystretch}{1.3}
\caption{Extensive Margin of Agricultural Land Use, 1974-2010 \label{tab:extensive}}
\sbox{\tabbox}{%
\begin{tabular}{l*{4}{c}}
\hline\hline
                    &\multicolumn{2}{c}{Sensitive}              &\multicolumn{2}{c}{Insensitive}            \\\cmidrule(lr){2-3}\cmidrule(lr){4-5}
                    &\multicolumn{1}{c}{(1)}         &\multicolumn{1}{c}{(2)}         &\multicolumn{1}{c}{(3)}         &\multicolumn{1}{c}{(4)}         \\
                    &         LPM         &      Probit         &         LPM         &      Probit         \\
\hline
$\Delta Log Travel Time$&       0.216\sym{**} &       0.230\sym{*}  &       0.076         &       0.134         \\
                    &     (0.101)         &     (0.131)         &     (0.065)         &     (0.138)         \\
Observations        &         528         &         394         &         528         &         383         \\
First-stage F stat  &      58.304         &                     &      58.304         &                     \\
\hline\hline
\end{tabular}
}
\usebox{\tabbox}
\begin{minipage}{\wd\tabbox}
    \vspace{0.1cm}
    {\footnotesize \sloppy \textit{Notes: } Each observation is a micro-region. Log change in travel time on roads is instrumented by the change in travel times on the predicted road network. The predicted road network is created by geo-referencing the PNV National Transportation Plan in 1973. Dependent variables are defined as a dummy equal to 1 if the micro-region recorded no area in the category in 1974 but positive area in 2010. The reported coefficients for Probit Regressions are the directly interpretable marginal effect. ***,**,* imply significance at the 1\%, 5\%, and 10\% levels, respectively.}
\end{minipage}
\endgroup
\end{table}
